\chapter[1944 --- Gasser]{1944: Joseph Erlanger and Herbert Spencer Gasser}
\label{ch:1944_Erlanger}

\section*{Main Contribution}

\textit{Demonstrated that individual nerve fibers carry highly specific information about different sensory modalities, explaining how the nervous system encodes sensory information.}

\section{Official Citation}

for their discoveries relating to the highly differentiated functions of single nerve fibres

\section{Background and Historical Context}

\subsection{Joseph Erlanger}

Joseph Erlanger was born on January 5, 1874, in San Francisco, California, USA. He was the son of Jewish immigrants from Germany who had settled in California during the Gold Rush era. His father, a merchant, encouraged his son's interest in science, and young Joseph pursued his education with determination and enthusiasm. He attended the University of Wisconsin, where he earned his bachelor's degree in 1895, and then moved to Johns Hopkins University, where he received his medical degree in 1899.

After completing his medical training, Erlanger joined the faculty of Johns Hopkins University, where he worked under the mentorship of William Henry Welch, one of the founders of the Johns Hopkins School of Medicine. At Johns Hopkins, Erlanger became interested in the physiology of the cardiovascular system, and he made important contributions to the understanding of blood pressure regulation and the function of the heart. In 1906, he published a landmark textbook of physiology (with Gasser) that became the standard reference in the field for many years.

In 1910, Erlanger moved to Washington University in St. Louis, where he became professor of physiology and later chairman of the Department of Physiology. It was at Washington University that he would conduct the research that led to his Nobel Prize, in collaboration with his former student and prot\'{e}g\'{e}, Herbert Spencer Gasser.

\subsection{Herbert Spencer Gasser}

Herbert Spencer Gasser was born on August 5, 1888, in Platteville, Wisconsin, USA. He grew up in a small Midwestern town, the son of a pharmacist, and developed an early interest in biology and medicine. He attended the University of Wisconsin, where he earned his bachelor's degree in 1910 and his medical degree in 1915. During his time at Wisconsin, Gasser studied under the influence of several prominent physiologists, including S.\ J.\ Meltzer, who had introduced him to the experimental study of the nervous system.

After completing his medical degree, Gasser joined Erlanger's laboratory at Washington University, where the two began their collaboration on the electrophysiology of nerve fibers. Their partnership would prove to be one of the most productive in the history of physiology, resulting in a series of landmark publications that fundamentally changed our understanding of how nerve impulses are conducted.

Gasser served as a military physician during World War I, an experience that deepened his interest in the physiology of the nervous system. After the war, he returned to Washington University and resumed his collaboration with Erlanger. In 1931, Gasser moved to Cornell University Medical College in New York, where he became professor of pharmacology and later dean of the medical school.

The scientific context of the early twentieth century was one of intense interest in the nervous system. Santiago Ram\'{o}n y Cajal had established the neuron doctrine, demonstrating that the nervous system is composed of individual nerve cells, and Charles Sherrington had described the concept of the synapse. However, the mechanisms by which nerve impulses travel along nerve fibers were poorly understood. The prevailing view, influenced by the work of Hermann von Helmholtz, held that nerve impulses are electrical in nature, but the details of this process---including the relationship between nerve fiber diameter and conduction velocity---were unknown. It was within this intellectual framework that Erlanger and Gasser conducted their pioneering research.

\section{The Discovery/Contribution}

Joseph Erlanger and Herbert Spencer Gasser received the Nobel Prize in Physiology or Medicine ``for their discoveries relating to the highly differentiated functions of single nerve fibres.'' Their work demonstrated that nerve fibers differ markedly in their diameter, myelination, and electrical properties, and that these differences are directly related to the speed and nature of nerve impulse conduction. This discovery fundamentally changed our understanding of the nervous system and laid the groundwork for modern neurophysiology.

\subsection{The Development of the Cathode Ray Oscilloscope Technique}

The key technical innovation that made Erlanger and Gasser's work possible was the adaptation of the cathode ray oscilloscope for the study of nerve impulses. Before their work, the electrical activity of nerve fibers could be detected using relatively crude galvanometers, but these instruments were too slow to capture the rapid changes in electrical potential that occur during nerve impulse conduction. Erlanger and Gasser recognized that the cathode ray oscilloscope---an instrument developed for use in television and radar---could provide the temporal resolution necessary to analyze nerve impulses in detail.

Working with a physicist, they developed an amplification system capable of detecting the tiny electrical signals generated by nerve fibers and displaying them on the oscilloscope screen. This technology allowed them, for the first time, to observe the shape, duration, and velocity of nerve impulses with notable precision. The technique represented a major advance in experimental physiology and opened up entirely new possibilities for the study of the nervous system.

\subsection{Relationship Between Fiber Diameter and Conduction Velocity}

Using their oscilloscope technique, Erlanger and Gasser analyzed the electrical activity of individual nerve fibers in the sciatic nerve of the frog and, later, in mammalian nerves. They discovered that the nerve impulses recorded from a mixed nerve---one containing many different types of fibers---consisted of multiple overlapping waves, each with a different velocity. By systematically analyzing these waves, they were able to correlate the velocity of nerve impulse conduction with the diameter of the nerve fiber.

Their findings were striking: larger diameter fibers conducted impulses more rapidly than smaller ones. Specifically, they found a linear relationship between fiber diameter and conduction velocity in myelinated fibers, with a conduction velocity-to-diameter ratio of approximately 6 meters per second per micrometer of diameter. This relationship---now known as the Erlanger-Gasser relationship---was one of the major quantitative discoveries in neurophysiology.

\subsection{Classification of Nerve Fibers}

Based on their observations of conduction velocity, Erlanger and Gasser developed the first systematic classification of nerve fibers. They identified three major groups, which they designated A, B, and C fibers:

\begin{itemize}
\item \textbf{A fibers} are large, heavily myelinated fibers with the fastest conduction velocity (12--120 meters per second). They are further subdivided into A$\alpha$, A$\beta$, A$\gamma$, and A$\delta$ subgroups based on their conduction velocities. A$\alpha$ fibers carry motor commands to skeletal muscles and sensory information from proprioceptors (muscle spindles and joint receptors), while A$\delta$ fibers carry signals for sharp, localized pain and temperature sensation.
\item \textbf{B fibers} are smaller, thinly myelinated fibers with intermediate conduction velocity (3--15 meters per second). They are found in the preganglionic fibers of the autonomic nervous system.
\item \textbf{C fibers} are the smallest and are unmyelinated, with the slowest conduction velocity (0.5--2 meters per second). They carry signals for dull, diffuse pain, temperature sensation, and autonomic functions.
\end{itemize}

This classification scheme, which has been refined but never fundamentally superseded, remains in use today and is taught in every medical school and neuroscience program around the world.

\subsection{The Role of Myelin in Nerve Conduction}

Erlanger and Gasser's work also provided the first experimental evidence for the role of the myelin sheath in determining nerve impulse conduction velocity. Myelin is a fatty insulating material that surrounds many nerve fibers, formed by specialized glial cells (Schwann cells in the peripheral nervous system and oligodendrocytes in the central nervous system). Erlanger and Gasser demonstrated that myelinated fibers conduct impulses more rapidly than unmyelinated fibers of comparable diameter, and they proposed that the myelin sheath enables a more efficient form of impulse conduction---saltatory conduction---in which the nerve impulse jumps from one node of Ranvier (gaps in the myelin sheath) to the next, rather than propagating continuously along the fiber.

This insight was ahead of its time and would not be fully validated until the development of intracellular recording techniques in the 1950s and 1960s. However, the hypothesis that myelin enables saltatory conduction was one of the major conceptual advances in neurophysiology and has since been confirmed by multiple lines of evidence.

\section{Impact on Modern Medicine}

The discoveries of Erlanger and Gasser have had significant implications for modern medicine, influencing our understanding and treatment of neurological diseases, the development of diagnostic tools, and the design of anesthetic techniques.

\subsection{Neurological Diagnosis}

The classification of nerve fibers by Erlanger and Gasser provides the foundation for modern neurophysiological diagnosis. Electromyography (EMG) and nerve conduction studies (NCS), which are among the most commonly used diagnostic tests in neurology, are based directly on the principles that Erlanger and Gasser established. By measuring the conduction velocity of nerve impulses along specific nerve fibers, clinicians can determine whether nerve damage has occurred, estimate the severity of the damage, and distinguish between different types of nerve disorders (e.g., demyelinating vs. axonal neuropathy).

The Erlanger-Gasser classification is also essential for understanding the effects of diseases that selectively affect specific nerve fiber types. Multiple sclerosis, for example, is a demyelinating disease that preferentially affects large, myelinated A fibers, leading to sensory loss, motor weakness, and impaired coordination. Conversely, certain peripheral neuropathies---such as those caused by diabetes or chemotherapy---preferentially affect small, unmyelinated C fibers, leading to pain, burning sensations, and loss of temperature sensation.

\subsection{Pain Management}

The distinction between different types of pain fibers---A$\delta$ fibers carrying sharp, acute pain and C fibers carrying dull, chronic pain---has been fundamental to the development of pain management strategies. The gate control theory of pain, proposed by Melzack and Wall in 1965, was based in part on the Erlanger-Gasser classification and proposed that the perception of pain can be modulated by the interaction between different types of nerve fibers. This theory has led to the development of transcutaneous electrical nerve stimulation (TENS) and other neuromodulation techniques for pain management.

Understanding the different types of pain fibers has also informed the development of pharmacological approaches to pain management. Local anesthetics, for example, preferentially block small, unmyelinated fibers at low concentrations, which explains why they can eliminate pain without completely abolishing motor function. This selectivity is the basis for epidural anesthesia, which is widely used in obstetrics and surgery.

\subsection{Anesthesiology}

The work of Erlanger and Gasser has been essential for the development of modern anesthesiology. Anesthesiologists must have a thorough understanding of the different types of nerve fibers and their functions in order to select appropriate anesthetic agents and techniques for different surgical procedures. The Erlanger-Gasser classification provides the framework for understanding why different anesthetic agents have different effects on different types of nerve fibers, and why regional anesthetic techniques can selectively block specific sensory or motor functions.

\subsection{Cardiac Electrophysiology}

The principles of nerve impulse conduction elucidated by Erlanger and Gasser have also influenced our understanding of cardiac electrophysiology. The heart's conduction system, which coordinates the rhythmic contraction of the cardiac muscle, shares many features with the nerve fibers that Erlanger and Gasser studied. Understanding the relationship between fiber diameter, myelination, and conduction velocity has been essential for the development of techniques such as cardiac ablation, which is used to treat arrhythmias (abnormal heart rhythms).

\section{Legacy}

The legacies of Joseph Erlanger and Herbert Spencer Gasser extend far beyond their specific discoveries to encompass the broader fields of neurophysiology and clinical neurology. Their classification of nerve fibers remains a cornerstone of neurophysiology, and the experimental techniques they pioneered---particularly the use of the cathode ray oscilloscope for the study of biological electrical signals---opened up entirely new possibilities for the investigation of the nervous system.

Erlanger continued his research at Washington University after receiving the Nobel Prize, making additional contributions to the understanding of cardiovascular physiology. He retired in 1944 but remained active in scientific circles until his death on December 5, 1965, at the age of 91. Gasser served as director of the National Heart Institute (now the National Heart, Lung, and Blood Institute) from 1950 to 1953, where he played a key role in establishing the institute's research programs. He continued his research on pain physiology until his retirement and died on May 11, 1963, at the age of 74.

The work of Erlanger and Gasser continues to influence modern neuroscience. Their classification of nerve fibers has been refined and extended by subsequent research, but the fundamental principles they established---that nerve fibers differ in their structural and functional properties, and that these differences are directly related to their physiological roles---remain as valid today as they were when first proposed in the 1920s and 1930s. Their discoveries stand as a testament to the power of careful experimentation and innovative technology in advancing our understanding of the nervous system.


\section{Modern Developments}

Erlanger and Gasser's work on nerve fiber function has informed modern neurology. Nerve conduction studies diagnose peripheral neuropathies and radiculopathies. Understanding of nerve fiber conduction velocity has led to treatments for demyelinating diseases (multiple sclerosis, Guillain-Barré syndrome). Peripheral nerve stimulation is used for chronic pain management. Brain-computer interfaces decode neural signals from peripheral nerves, enabling prosthetic control.


\section{Bibliography}

\begin{thebibliography}{9}
\bibitem{nobel-1944}
Nobel Prize Outreach, ``The Nobel Prize in Physiology or Medicine 1944,''\emph{NobelPrize.org}.\\
\url{https://www.nobelprize.org/prizes/medicine/1944/summary/}

\bibitem{lecture-1944}
Joseph Erlanger, ``Nobel Lecture,''\emph{NobelPrize.org}. Winner-authored primary source; downloadable from the official Nobel Prize archive.\\
\url{https://www.nobelprize.org/prizes/medicine/1944/erlanger/lecture/}
\end{thebibliography}
